\documentclass[10pt]{article}

\usepackage[a4paper,margin=20mm]{geometry}
\usepackage[T1]{fontenc}
\usepackage{lmodern}
\usepackage[hidelinks]{hyperref}
\usepackage{microtype}
\usepackage{parskip}

\hypersetup{
  pdftitle={Cover letter - Layered representation diagnostics for short metagenomic reads},
  pdfauthor={Ruixiang Mei}
}

\pagestyle{empty}
\setlength{\parindent}{0pt}
\setlength{\parskip}{5pt}

\begin{document}

\begin{flushright}
Ruixiang Mei\\
School of Data Science\\
The Chinese University of Hong Kong, Shenzhen\\
2001 Longxiang Road, Longgang District\\
Shenzhen 518172, Guangdong, China\\
\href{mailto:121090416@link.cuhk.edu.cn}{121090416@link.cuhk.edu.cn}\\
\href{https://orcid.org/0009-0003-2128-0726}{ORCID: 0009-0003-2128-0726}\\[3pt]
\today
\end{flushright}

Editors\\
\textit{NAR Genomics and Bioinformatics}

Dear Editors,

Please consider the manuscript entitled ``Layered representation diagnostics for short metagenomic reads'' for publication as a Standard Paper in \textit{NAR Genomics and Bioinformatics}.

Short-read representations are commonly judged only through downstream predictive accuracy, which does not reveal whether local composition, biochemical summaries or coarse positional information remain accessible after compression and perturbation. This manuscript introduces a paired representation-diagnostics framework and a 222-dimensional, training-free representation, CK4P-MSP, that combines reverse-complement canonical 4-mer composition with global biochemical properties and multi-scale positional property pooling. Its main methodological contribution is a block-decomposable audit construction in which these information channels can be examined under one controlled perturbation protocol.

The evidence combines a seven-group K/P/MSP ablation, shared-template length and perturbation sweeps, matched counterfactual controls, dimension-matched high-k vectors, PCA/SVD reductions, historical biochemical descriptors and full-position upper bounds. The analyses associate the global property block primarily with reduced perturbation drift, the multi-scale positional block with grouped local-change readability, and the canonical k-mer block with composition-linked retrieval. PseKNC and full-position encodings define lower-drift and higher-positional-readout boundaries, respectively. ART and public CAMI resources then provide simulator-contextualized stability, coarse fixed-head label retention and an anonymous-read composition-shift probe beyond the controlled whole-genome-sequence grid.

We believe the manuscript is well suited to the journal because it addresses a general genomics and bioinformatics problem: how to evaluate what compact sequence representations retain. It does so through quantitative comparisons, openly specified representation contracts and a reproducible workflow. Source code, processed source tables, figure inputs and a manuscript-to-script map are maintained in an MIT-licensed versioned repository at \url{https://github.com/FairYJmrx/ultrashort-dna-representation-diagnostics}. The frozen \texttt{v1.0.0} release is archived at \url{https://doi.org/10.5281/zenodo.21792508}.

This manuscript is original, has not been published and is not under consideration by another journal. No preprint has been posted as of the date of this letter. The study used only public reference and benchmark resources and computationally simulated sequences; no human participants, patient-level data or identifiable private information were involved. The author declares no competing interests and received no specific funding for this work.

Generative-AI tools (OpenAI ChatGPT and Codex, Google Gemini, and Doubao) were used primarily as assistive review and quality-control tools for language, terminology, internal consistency, references, formatting and code, including support for preparing and checking plotting scripts. No AI-generated output was used as experimental data or treated as scientific evidence. Scientific decisions, interpretation and final wording remained under the author's control. The author independently reviewed and verified the manuscript, code, calculations, references and figures.

Thank you for considering this manuscript. I would be grateful for the opportunity to have it evaluated by the journal.

Sincerely,\\[8pt]
Ruixiang Mei\\
Corresponding author

\end{document}
